\scanpage{280}
\mspage{10}

\uline{Proposition} 2.\unsure{4|7}. Soit $W$ un localiseur fondamental. Alors les produits finis existent dans $\mathrm{Hot}_W$, et le foncteur $\mathrm{hot}_W:\mathrm{Cat}\unsure{{-}|\longrightarrow}\mathrm{Hot}_W\overset{\mathrm{def}}{=}(\mathrm{Cat})W^{-1}$ y commute. Si $W$ satisfait à \uline{Loc} \deleted{\unsure{3}}$4a$ (resp. \uline{Loc} $4a_0$), alors les sommes quelconques \inserted{(resp. finies)} existent dans $\mathrm{Hot}_W$, et le foncteur $\mathrm{hot}_W$ y commute. De plus, dans ce cas les sommes sont distributives p.r. aux produits (vu qu'il en est ainsi dans Cat).

La démonstration résulte simplement du résultat général :

\uline{Lemme} 2.5. Soit $(M,W)$ une catégorie de modèles ($W$ partie quelconque de $\mathrm{Fl}\,M$). Supposons \inserted{\unsure{que}} \deleted{toute famille \unsure{d'éléments}}\inserted{les produits \deleted{finis}} (resp. les produits \deleted{quelconques}\inserted{finis}) \inserted{(resp. les produits binaires)} existent dans $M$, et que de plus pour \deleted{\unsure{la|toute}} famille $(f_i)_{i\in I}$ de flèches de $M$, \deleted{telle} (resp tte famille finie, resp. \deleted{pour deux flèches}\inserted{famille binaire}), \deleted{\unsure{tels}} $f_i\in W$ pour tt $i\in I$ implique $f=\prod f_i\in W$,

Alors les produits \unsure{au type|du type|au degré} envisagé existent dans $\mathrm{Ho}_W=MW^{-1}$, et le foncteur \deleted{$\unsure{\mathrm{loc}_W}$} $\unsure{r_W|i_W}:M{-}MW^{-1}$ y commute. Enfin, si $M$ admet un objet final $e$, alors $e$ est aussi objet final dans $\mathrm{Ho}_W$.

% archive sheet number: 279
